\section{Hidden affine consequences: quartics in 278 and 279 integers}
\label{sec:row-exposure}

The new direct-calculus contribution is a reduction of the incoming
282-coordinate quartic to 279 coordinates without a support or height
penalty, and to 278 coordinates at a stated support tradeoff.  The mechanism
is not a new number-theoretic assumption: elementary changes of residual
basis expose affine identities that the original gate presentation hides.
The corresponding 287- and 309-coordinate parents also lose three
coordinates.  Every coordinate remains an unrestricted integer.

\begin{center}\small
\begin{tabular}{lrrrr}
\hline
Equation & Variables & Rows & Monomials & Height\\
\hline
Incoming dimension parent &282&192&5491&128446\\
Incoming grouped parent &287&197&4313&128446\\
Incoming sparse parent &309&219&3922&128446\\
\hline
New dimension tradeoff &278&188&12884&12227\\
New balanced &279&189&5483&7200\\
New sparse 279 &279&189&5202&12662\\
New sparse 284 &284&194&4270&10864\\
New sparse 306 &306&216&3870&14512\\
\hline
\end{tabular}
\end{center}

All eight displayed equations have exact degree four.  The two
279-coordinate equations improve every varying metric of the matched
282-coordinate parent; the 284- and 306-coordinate equations do likewise
for their matched parents.  The 278-coordinate point sacrifices support
for one fewer variable.  These are literal finite endpoints, not
global-minimality or literature-priority claims.

\subsection{Equivalence lemmas}

\paragraph{Residual basis.}
If $i\ne j$ and $k\in\mathbb Z$, replacing $R_i$ by $R_i+kR_j$
leaves the common zero set unchanged: one implication is substitution,
and the other follows by subtracting $kR_j$ after using $R_j=0$.
The inverse uses $-k$, so each such operation is unimodular.  The
sum-of-squares polynomial generally changes; its zero set does not.

\paragraph{Integral coordinate chart.}
The substitution $x_t=z_t+L(z)$, with $L$ integral affine and independent
of $z_t$, is a bijection of integer lattices.  Its inverse is
$z_t=x_t-L(x)$.  A sequence is inverted in reverse order.  Quadratic
degree is preserved.

\paragraph{Unit affine projection.}
If one residual is $\epsilon x_t+L=0$, with $\epsilon\in\{1,-1\}$
and $L$ integral affine and independent of $x_t$, substitute
$x_t=-\epsilon L$ into every other residual and delete this row and
coordinate.  Restriction maps old zeros to new zeros; the displayed
integral formula gives the unique inverse extension.  Repeated extensions
are evaluated in reverse projection order.  Degree remains at most two.

Since $\sum_i R_i^2=0$ over $\mathbb Z$ exactly when every $R_i=0$,
these three facts prove the corresponding equivalence for a single quartic.
No sign guard, divisibility condition or cofinal choice is introduced.

\subsection{Credited preconditioning and the new identities}

Six previously published-in-repository coordinate shears are transferred
to the three incoming parents.  This preconditioning is credited prior
work, not the novelty claim.  With zero-based coordinate indices, the
successive old-coordinate substitutions are
\[
\begin{aligned}
x_{21}&\leftarrow x_{21}-2x_{10}-2x_{11},&
x_{23}&\leftarrow x_{23}-8x_{22},\\
x_{38}&\leftarrow x_{38}-8x_{37},&
x_{35}&\leftarrow x_{35}-4x_{10}-5x_{12}-2x_{25}-x_{27},\\
x_{27}&\leftarrow x_{27}-x_{10}-x_{12},&
x_{25}&\leftarrow x_{25}-x_{11}-x_{12}.
\end{aligned}
\]
They are executed in the displayed order, not simultaneously.  The
upstream certificates, exact source hashes and commit are retained.
Each forward and inverse map is independently checked here; no semantic
acceptance claim for the upstream V-R1 endpoint is needed.

Let $Q_i$ denote the residuals after this preconditioning.  Exact
coefficient collection in each of the three parents gives
\[
\begin{aligned}
Q_{21}-Q_6
 &=4+4x_{10}+8x_{11}+4x_{21}-x_{22}+x_{37},\\
Q_{73}-Q_{71}-4Q_{70}-8Q_{68}
 &=12-6x_3-8x_{84}-4x_{85}-x_{86}+x_{89},\\
Q_{122}-Q_{120}
 &=3-2x_{107}+2x_{151}-x_{167}+x_{169}.
\end{aligned}
\]
These are polynomial identities, not consequences inferred from sample
solutions.  Replace rows 21, 73 and 122 by the left sides, using five
elementary row shears.  The donor rows are unchanged.  The unit-affine
projection lemma then gives
\[
\begin{aligned}
x_{37}&=x_{22}-4-4x_{10}-8x_{11}-4x_{21},\\
x_{89}&=x_{86}-12+6x_3+8x_{84}+4x_{85},\\
x_{169}&=x_{167}-3+2x_{107}-2x_{151}.
\end{aligned}
\]
The deleted coordinates are named \texttt{add.145}, \texttt{mul.346}
and \texttt{base4.T}.  This proves the drops
$(282,192)\to(279,189)$, $(287,197)\to(284,194)$ and
$(309,219)\to(306,216)$ in coordinates and rows.

\subsection{One further projection and the 278-coordinate point}

The already affine row $Q_{90}$ has unit coordinate $x_{105}$,
named \texttt{add.380}.  It is not counted as a newly discovered affine
identity.  Its unique restoration is
\[
\begin{aligned}
x_{105}={}&x_{10}+x_{11}+4x_{21}+2x_{24}+4x_{26}+2x_{29}
            +2x_{31}+4x_{32}\\
          &+2x_{34}+2x_{36}+4x_{39}+2x_{40}+2x_{43}
            +2x_{47}+4x_{68}+2x_{69}.
\end{aligned}
\]
Applying this additional projection to the dimension parent gives 278
coordinates and 188 rows.  The dense replacement explains the larger
collected support.  No new semantic lemma is needed.

\subsection{The fixed quartics and their logical meaning}

Each endpoint is defined exactly by
\[
 P(\mathbf z)=\sum_{i=0}^{m-1}r_i(\mathbf z)^2,
 \qquad \mathbf z\in\mathbb Z^n,
\]
where its complete list of integral quadratic $r_i$ is in
\texttt{equation.json}; complete collected coefficients are in
\texttt{expanded\_polynomial.json}.  The package also supplies every row
as human-readable text and TeX, the full coordinate-order table, and a
complete affine map restoring the original parent coordinates.

Following the projections, a frozen list of elementary residual row
shears selects each stated support/height profile.  The lists are fully
explicit in the endpoint certificates.  Their proof is the residual-basis
lemma applied successively, not the search heuristic that found them.
The deterministic builder runs no optimizer.  Exact expansion confirms
degree four and that every declared coordinate occurs.

\paragraph{Theorem.}
Each supplied new equation has an integer zero if and only if its
matched incoming parent has an integer zero.  Hence, composing with the
parent's ordinary C222/ZF/ZFC refutation calibration, ZFC proves that
integer unsolvability of each new equation is equivalent to
$\operatorname{Con}(\mathrm{ZFC})$.

\paragraph{Proof.}
For a new zero, undo the final row shears to recover all projected
residual equations.  Restore the deleted coordinates in reverse order;
the deleted equations vanish identically and all remaining equations
hold.  Undo the affine-exposure row shears, then evaluate the coordinate
charts in reverse execution order.  This produces an old zero.  In the
other direction, use the inverse coordinate bijection, the exact affine
equalities, restriction to retained coordinates, and the forward row
shears.  Every step is integral and reversible on zeros.  The final
logical equivalence is exactly the inherited parent calibration.
This argument does not prove consistency or claim a new end-to-end
kernel-checked Lean theorem.

\subsection{Verification and the limit of this particular method}

An independent standard-library checker imports no producer.  It checks
coordinate inverses, elementary row inverses, unit coefficients, the
identical zero of every deleted row, absence of deleted coordinates,
the exact coordinate renumbering, every collected coefficient and metric,
and coordinate liveness.  It also reconstructs every original coordinate
as an affine polynomial in the final ones and checks every original
residual pullback as a polynomial identity.  Five corrupted-artifact
controls are rejected.  Normal and \texttt{python -O} runs pass, including
fresh byte-identical rebuilds.  The printed residual text and TeX are
parsed back to their complete coefficient lists.

For the 278-coordinate endpoint, the quadratic parts of the 188 residuals
have rank 188 modulo 1000003.  Thus they also have rank 188 over
$\mathbb Q$: a nonzero modular minor cannot have zero integer
determinant.  No nonzero rational linear combination of this final
residual basis is affine.  The parent's corresponding rank is 188, and
the four displayed independent affine combinations account for its
four-dimensional kernel.  This finite obstruction concerns this row
span only, not nonlinear ideal consequences or global variable
minimality.  The optional rank replay is supplied separately.
